% !PRINTED 352
\weberpage{352}

%% [running head, not body text -- commented out 2026-08-07] Page furniture (running
%% head, folio, signature) is excluded from this edition by standing policy.
%% \noindent 352\hfill Zwölfter Abschnitt.\hfill \S\ 96.

Wir setzen, wenn $\Delta$ die Grundzahl des Körpers $\Omega$ ist, je nachdem $\Delta$ gerade oder ungerade ist:
\weq{1}{
\begin{aligned}
\theta_0 &= \tfrac{1}{2}\sqrt{\Delta}, & &\text{oder}\\
\theta_0 &= \frac{-1+\sqrt{\Delta}}{2}\,.
\end{aligned}
}

Dann ist nach \S\ 90 jede ganze Zahl des Körpers $\Omega$ in der Form
\weq{2}{
\omega = x + y\theta_0
}
darstellbar, worin $x$ und $y$ ganze Zahlen sind, und diese Zahl ist dann und nur dann nach dem Modul $Q$ mit einer rationalen Zahl $a$ kongruent, wenn $y$ durch $Q$ teilbar ist, weil dann
\[
\frac{x - a + y\theta_0}{Q}
\]
eine ganze Zahl sein muß. Demnach sind alle Zahlen der Ordnung $[Q]$ in der Form enthalten:
\[
\omega = x + yQ\theta_0,
\]
oder in den beiden in (1) unterschiedenen Fällen, wenn
\weq{3}{
D = Q^2\Delta
}
gesetzt wird:
\weq{4}{
\begin{aligned}
&x - y\,\frac{\sqrt{D}}{2}, & &\Delta \text{ gerade},\\
&x - \frac{Q}{2}y + y\,\frac{\sqrt{D}}{2}, & &\Delta \text{ ungerade}, D \text{ gerade},\\
&x - \frac{Q-1}{2}y + y\,\frac{-1+\sqrt{D}}{2}, & &D \text{ ungerade}.
\end{aligned}
}

Setzen wir also
\weq{5}{
\begin{aligned}
\theta &= \frac{\sqrt{D}}{2}, & D &\equiv 0 \pmod{4}\\
\theta &= \frac{-1+\sqrt{D}}{2}, & D &\equiv 1 \pmod{4}
\end{aligned}
}
und ersetzen in den beiden letzten Formeln (4) die willkürlich ganzen Zahlen
\[
x - \frac{Q}{2}y, \quad x - \frac{Q-1}{2}y
\]
durch $x$, so erhalten wir das Resultat:
%CONT>>
